\documentclass[10pt, aspectratio=169]{beamer}
\usepackage{../beamer_theme_408}

\title{408考研零基础 --- 计算机组成原理}
\subtitle{2-15 流水线技术}
\author{}
\date{}

\begin{document}


\begin{frame}{本节目标}
    \begin{enumerate}
        \item 理解指令流水线的基本概念
        \item 掌握流水线的三大性能指标
        \item 区分三种流水线冒险
        \item 了解超标量与多发射技术
    \end{enumerate}
\end{frame}

\begin{frame}{流水线的基本概念}
    \begin{block}{思想来源 \ --- \ 工厂装配线}
        传统：一个工人完成全部工序 $\rightarrow$ 效率低 \\
        流水线：各工人专攻一道工序 $\rightarrow$ 并行度高
    \end{block}
    \vspace{0.3em}
    \textbf{指令流水线}：将指令执行过程分成多个阶段，各阶段重叠并行
    \vspace{0.3em}
    \begin{center}
        \begin{tabular}{|c|c|c|c|c|}
            \hline
            时钟 & IF & ID & EX & WB \\
            \hline
            1 & 指令1 &  &  &  \\
            2 & 指令2 & 指令1 &  &  \\
            3 & 指令3 & 指令2 & 指令1 &  \\
            4 & 指令4 & 指令3 & 指令2 & 指令1 \\
            \hline
        \end{tabular}
    \end{center}
    \vspace{0.3em}
    \begin{itemize}
        \item IF（取指）、ID（译码）、EX（执行）、WB（写回）
    \end{itemize}
\end{frame}

\begin{frame}[fragile]{流水线的五个经典阶段}
    经典五级流水线（RISC）：
    \vspace{0.3em}
    \begin{enumerate}
        \item \textbf{IF}（Instruction Fetch）--- 取指令
            \begin{itemize}
                \item PC $\rightarrow$ 指令存储器
            \end{itemize}
        \item \textbf{ID}（Instruction Decode）--- 译码
            \begin{itemize}
                \item 读寄存器堆，译码
            \end{itemize}
        \item \textbf{EX}（Execute）--- 执行
            \begin{itemize}
                \item ALU运算或地址计算
            \end{itemize}
        \item \textbf{MEM}（Memory Access）--- 访存
            \begin{itemize}
                \item 读/写数据存储器
            \end{itemize}
        \item \textbf{WB}（Write Back）--- 写回
            \begin{itemize}
                \item 将结果写回寄存器堆
            \end{itemize}
    \end{enumerate}
\end{frame}

\begin{frame}{流水线的性能指标 --- 吞吐率}
    \textbf{吞吐率（TP）}：单位时间内流水线完成的指令数量
    \vspace{0.3em}
    \begin{itemize}
        \item 理想情况：
            \begin{itemize}
                \item $n$ 条指令，$k$ 级流水线，时钟周期 $t$
                \item 总时间 $T = (k + n - 1) \times t$
                \item $\text{TP} = \dfrac{n}{T} = \dfrac{n}{(k + n - 1) \times t}$
            \end{itemize}
        \item 最大吞吐率：
            \begin{itemize}
                \item $n \rightarrow \infty$ 时，$\text{TP}_{\max} = \dfrac{1}{t}$
            \end{itemize}
    \end{itemize}
\end{frame}

\begin{frame}{流水线的性能指标 --- 加速比与效率}
    \textbf{加速比（S）}
    \begin{itemize}
        \item 流水线执行时间 vs 非流水线执行时间
        \item $S = \dfrac{T_{\text{非流水}}}{T_{\text{流水}}}$
        \item 理想情况 $S_{\max} = k$（$k$ 为流水线段数）
    \end{itemize}
    \vspace{0.3em}
    \textbf{效率（E）}
    \begin{itemize}
        \item 流水线中各功能段的利用率
        \item $E = \dfrac{n \times k}{(k + n - 1) \times k} = \dfrac{n}{k + n - 1}$
        \item $n \rightarrow \infty$ 时，$E \rightarrow 1$
    \end{itemize}
    \vspace{0.3em}
    \begin{block}{三个指标的关系}
        $\text{TP} = \dfrac{E}{t}$，$S = k \times E$
    \end{block}
\end{frame}

\begin{frame}{流水线冒险 --- 结构冒险}
    \textbf{结构冒险（Structural Hazard）}
    \begin{itemize}
        \item 硬件资源冲突导致无法同时完成指令
        \item 例：指令存储器与数据存储器共用同一端口
    \end{itemize}
    \vspace{0.3em}
    \textbf{解决方法}
    \begin{itemize}
        \item 分离指令Cache和数据Cache（哈佛结构）
        \item 插入气泡（Stall / Bubble）
        \item 资源重复（Resource Replication）
    \end{itemize}
    \vspace{0.3em}
    \begin{center}
        \begin{tabular}{|c|c|c|c|c|}
            \hline
            时钟 & IF & ID & EX & MEM \\
            \hline
            3 & 指令3 & 指令2 & 指令1 &  \\
            4 & \textcolor{red}{气泡} & 指令3 & 指令2 & 指令1 \\
            \hline
        \end{tabular}
    \end{center}
\end{frame}

\begin{frame}{流水线冒险 --- 数据冒险}
    \textbf{数据冒险（Data Hazard）}
    \begin{itemize}
        \item 后续指令依赖前面指令的运算结果
        \item 类型：RAW（先写后读）、WAR（先读后写）、WAW（先写后写）
    \end{itemize}
    \vspace{0.3em}
    \textbf{例（RAW）}：
    \begin{itemize}
        \item \texttt{ADD R1, R2, R3}
        \item \texttt{SUB R4, R1, R5} \quad （依赖R1结果）
    \end{itemize}
    \vspace{0.3em}
    \textbf{解决方法}
    \begin{itemize}
        \item 数据转发（Forwarding / Bypassing）
        \item 插入气泡（软件插入NOP或硬件阻塞）
        \item 编译器重排指令
    \end{itemize}
\end{frame}

\begin{frame}{流水线冒险 --- 控制冒险}
    \textbf{控制冒险（Control Hazard）}
    \begin{itemize}
        \item 遇到转移指令（JMP、BRZ等）时，流水线不确定下一条指令地址
    \end{itemize}
    \vspace{0.3em}
    \textbf{解决方法}
    \begin{enumerate}
        \item \textbf{分支预测}：
            \begin{itemize}
                \item 静态预测：总预测跳转 / 不跳转
                \item 动态预测：根据历史记录预测
            \end{itemize}
        \item \textbf{延迟槽}：
            \begin{itemize}
                \item 在转移指令后插入一条总是执行的指令
            \end{itemize}
        \item \textbf{清空流水线}：
            \begin{itemize}
                \item 预测错误时冲刷流水线，性能损失大
            \end{itemize}
    \end{enumerate}
\end{frame}

\begin{frame}{超标量与多发射技术}
    \textbf{超标量（Superscalar）}
    \begin{itemize}
        \item 每个时钟周期发射多条指令
        \item 需多个功能单元并行工作
        \item 例：Intel Core 系列
    \end{itemize}
    \vspace{0.3em}
    \textbf{多发射（Multiple Issue）}
    \begin{itemize}
        \item 静态多发射：编译器决定哪些指令并行发射
        \item 动态多发射：硬件动态调度指令发射
    \end{itemize}
    \vspace{0.3em}
    \begin{center}
        \begin{tabular}{c|c|c|c}
            \textbf{类型} & \textbf{IPC} & \textbf{调度方式} & \textbf{复杂度} \\
            \hline
            单发射 & 1 & 简单 & 低 \\
            超标量 & $>1$ & 动态 & 高 \\
            超长指令字（VLIW） & $>1$ & 静态 & 中 \\
        \end{tabular}
    \end{center}
\end{frame}

\begin{frame}{本节总结}
    \begin{enumerate}
        \item 流水线将指令分阶段并行，提高吞吐率
        \item 性能指标：吞吐率 TP、加速比 S、效率 E
        \item 三大冒险：结构冒险（资源冲突）、数据冒险（数据依赖）、控制冒险（转移指令）
        \item 超标量技术每个周期发射多条指令，进一步提升性能
    \end{enumerate}
\end{frame}

\end{document}
